
\section{\hard\ 从射影到固有假设：\trChow{}引理与\trGrothendieck{}凝聚定理}

本节中的各个主证明都标以双星，原因在于这些结果在我们其余的讨论中并非绝对必要，现在也许不值得去读。不过，只要知道\mbox{\trGrothendieck{}凝聚定理~18.9.1}（它推广了\mbox{定理~18.7.1(d)}）的陈述，你就能立刻把我们的许多关于射影概形与射影态射的论证移植到固有概形与固有态射上去，而\trChow{}引理则是一件多用途的工具，可以把结果从射影情形推广到一般的固有情形。

于是你有四种选择：$(i)$ 通读\trGrothendieck{}凝聚定理的证明；$(ii)$ 把证明浏览一遍，说服自己：只要你想懂，你就\emph{能}懂；$(iii)$ 跳过证明，径直假定它为真；或者 $(iv)$ 每当\trGrothendieck{}凝聚定理被调用时，就干脆退回到射影假设。人生充满了妥协，而我不会对你评头论足。

\textbf{18.9.1.\ \trGrothendieck{}凝聚定理}\textbf{。}\ ——\ \emph{设 $\pi\colon X\to Y$ 是\trLocally\trNoetherian{}概形之间的一个\trProper\trMorphism。则对 $X$ 上任意\trCoherentSheaf{} $\mathscr{F}$，$\mathrm{R}^i\pi_*\mathscr{F}$ 在 $Y$ 上是\trCoh{}的。}

情形 $i=0$ 已经提到过好几次了。

\textbf{18.9.A.（练习）}\quad 回顾一下：有限态射是仿射的（按定义），并且是固有的。利用\mbox{定理~18.9.1} 证明：若 $\pi\colon X\to Y$ 既固有又仿射，而 $Y$ 是\trNoetherian{}的，则 $\pi$ 是有限的。（提示：模仿把固有换成射影的那个较弱结果的证明，\mbox{推论~18.1.5}。）

\mbox{定理~18.9.1} 的证明需要两个精深的事实。第一个是\trLeray\trSpectralSequence（\mbox{定理~23.4.5}；由于\mbox{练习~23.5.F}，它适用于这里的情形）。设 $\pi\colon X\to Y$ 与 $\rho\colon Y\to Z$ 是\trQuasicompact\trSeparated\trMorphism。那么对 $X$ 上任意\trQCSheaf{} $\mathscr{F}$，存在一条\trSpectralSequence，其 $E_2$ 项由 $\mathrm{R}^p\rho_*(\mathrm{R}^q\pi_*\mathscr{F})$ 给出，并\trConverge{}到 $\mathrm{R}^{p+q}(\rho\circ\pi)_*\mathscr{F}$。由于在我们需要的情形（即 $Z$ 仿射的情形）里，这本来会是一道合理的（但有难度的）练习，所以我们用起来心安理得。不过我们稍后会证明它，现在就不证了。

我们还需要\trChow{}引理。

\textbf{18.9.2.\ \trChow{}引理}（\trProper{}＝\trProjective{}减去\trBirational\trProjective）\textbf{。}\ ——\ \emph{设 $\pi\colon X\to\Spec A$ 是一个\trProper\trMorphism，而 $A$ 是\trNoetherian{}的。则存在 $\mu\colon X'\to X$，它既是满射又是\trProjective{}的，使得 $\pi\circ\mu$ 也是\trProjective{}的，并且 $\mu$ 在 $X$ 的一个\trDense{}开子集上是同构：}
\[
\begin{tikzcd}
X' \ar[rr, "\mu"] \ar[dr, "\pi\circ\mu"'] && X \ar[dl, "\pi"] \\
& \Spec A. &
\end{tikzcd}
\]

特别地，若 $X$ 是一个\trProper{} $k$-\trVariety，则它容许一个来自某个\trProjective{} $k$-\trVariety{}的\trProjective\trBirational\trMorphism。

从射影情形到固有情形，许多结果的推广都要经由\trChow{}引理来实现。我们将在 \S18.9.4 中证明这个版本，并陈述\trChow{}引理的其他版本。假定这两个事实，我们现在就通过一系列练习来证明\mbox{\trGrothendieck{}凝聚定理~18.9.1}。

\textbf{18.9.3.}\ \extrahard\ \emph{\mbox{\trGrothendieck{}凝聚定理~18.9.1} 的证明，假定\mbox{\trChow{}引理~18.9.2}。}

这个问题在 $Y$ 上是局部的，所以不妨设 $Y$ 是仿射的，比如说 $Y=\Spec A$。我们对 $\operatorname{Supp}\mathscr{F}$ 施行诺特归纳，基础情形是 $\operatorname{Supp}\mathscr{F}=\varnothing$ 的情形，即 $\mathscr{F}=0$，这是显然的。于是取定 $\mathscr{F}$，并假设结论对一切支撑严格更小的\trCoherentSheaf{}来说已知。

\textbf{18.9.B.（练习）}\quad 证明：我们可以假设 $\operatorname{Supp}\mathscr{F}=X$。（提示：把 $X$ 换成 $\mathscr{F}$ 的\trSchemeTheoretic{}支撑，其定义见\mbox{练习~9.1.N}。）

现在我们援引\trChow{}引理，构造一个\trProjective\trMorphism{} $\mu\colon X'\to X$，它在 $X$ 的一个\trDense{}开子集 $U$ 上是同构（于是 $X\setminus U$ 是 $X$ 的一个\trClosed{}子集，严格小于 $X$），并且使得 $\pi\circ\mu\colon X'\to\Spec A$ 是\trProjective{}的。

于是 $\mathscr{G}=\mu^*\mathscr{F}$ 是 $X'$ 上的一个\trCoherentSheaf，$\mu_*\mathscr{G}$ 是 $X$ 上的一个\trCoherentSheaf（由射影情形，\mbox{定理~18.7.1(d)}），而\trAdjoint{}映射 $\mathscr{F}\to\mu_*\mathscr{G}=\mu_*\mu^*\mathscr{F}$ 在 $U$ 上是同构。它的\trKernel{} $\mathscr{E}$ 与\trCokernel{} $\mathscr{H}$ 是 $X$ 上的、支撑严格更小的\trCoherentSheaf：
\[
0\longrightarrow\mathscr{E}\longrightarrow\mathscr{F}\longrightarrow\mu_*\mathscr{G}\longrightarrow\mathscr{H}\longrightarrow 0.
\]

\textbf{18.9.C.（练习）}\quad 由归纳假设，$\mathscr{E}$ 与 $\mathscr{H}$ 的高阶\trPushforward{}都是\trCoh{}的。证明：若 $\mu_*\mathscr{G}$ 的高阶\trPushforward{}全是\trCoh{}的，则 $\mathscr{F}$ 的高阶\trPushforward{}也是\trCoh{}的。

于是我们化归为证明：对 $X'$ 上任意\trCoherentSheaf{} $\mathscr{G}$，$\mu_*\mathscr{G}$ 的高阶\trPushforward{}都是\trCoh{}的。

对于 $X'\xrightarrow{\ \mu\ }X\xrightarrow{\ \pi\ }\Spec A$ 的\trLeray\trSpectralSequence（\mbox{定理~23.4.5}），其 $E_2$ \trPage{}由 $\mathrm{R}^p\pi_*(\mathrm{R}^q\mu_*\mathscr{G})$ 给出，并\trConverge{}到 $\mathrm{R}^{p+q}(\pi\circ\mu)_*\mathscr{G}$。现在 $\mathrm{R}^q\mu_*\mathscr{G}$ 由\mbox{定理~18.7.1(d)} 知是\trCoh{}的。此外，由于 $\mu$ 在 $X$ 的一个\trDense{}开子集 $U$ 上是同构，$\mathrm{R}^q\mu_*\mathscr{G}$ 在 $U$ 上为零，从而支撑在 $U$ 的\trComplement{}上，此支撑严格小于 $X$。于是由我们的归纳假设，$\mathrm{R}^p\pi_*(\mathrm{R}^q\mu_*\mathscr{G})$ 对一切 $p$、一切 $q\ge 1$ 都是\trCoh{}的。$E_2$ \trPage{}上唯一可能不\trCoh{}的层都在 $q=0$ 这一行——恰好就是我们感兴趣的那些层。另外，把\mbox{定理~18.7.1(d)} 用于 $\pi\circ\mu$，可知 $\mathrm{R}^{p+q}(\pi\circ\mu)_*\mathscr{G}$ 是\trCoh{}的。

\textbf{18.9.D.（练习）}\quad 证明：对任意 $n\ge 2$ 与 $q>0$，$E_n^{p,q}$ 总是\trCoh{}的。证明：对一个给定的 $n\ge 2$，$E_n^{p,0}$ 是\trCoh{}的，当且仅当 $E_2^{p,0}$ 是\trCoh{}的。证明：$E_\infty^{p,q}$ 是\trCoh{}的，从而 $E_2^{p,0}$ 是\trCoh{}的，至此完成\mbox{\trGrothendieck{}凝聚定理~18.9.1} 的证明。\hfill$\square$

\textbf{18.9.4.}\ \extrahard\ \textbf{\mbox{\trChow{}引理~18.9.2} 的证明。}

我们要逐项动用 $\pi$ 之\trProperness{}假设的三个组成部分：有限型、分离、泛闭。其中用到分离性的部分尤其棘手。

由于 $X$ 是\trNoetherian{}的，它只有有限多个\trIrreducibleComponent{}。用仿射开集 $U_1,\dots,U_n$ 覆盖 $X$。我们可以假设每个 $U_i$ 都与每个\trIrreducibleComponent{} $Z$ 相遇。（若某个 $U_i$ 不与某个\trIrreducibleComponent{} $Z$ 相遇，则取 $Z\setminus\overline{X-Z}$ 中任一仿射开子集 $Z'$，并用 $U_i\cup Z'$ 代替 $U_i$。）于是 $U:=\cap_i U_i$ 是 $X$ 的一个\trDense{}开子集。由于每个 $U_i$ 在 $A$ 上都是有限型的，我们可以选取闭嵌入 $U_i\subset\mathbb{A}^{n_i}_A$。令 $\overline{U}_i$ 为 $U_i$ 在 $\mathbb{P}^{n_i}_A$ 中的（\trSchemeTheoretic{}）闭包。

现在我们有\trDiagonalMorphism{} $U\to X\times_A\prod\overline{U}_i$（其中乘积在 $\Spec A$ 上取），它是一个\trLocallyClosedEmbedding（闭嵌入 $U\hookrightarrow U^{n+1}$ 与开嵌入 $U^{n+1}\hookrightarrow X\times_A\prod\overline{U}_i$ 的复合）。令 $X'$ 为 $U$ 在 $X\times_A\prod\overline{U}_i$ 中的\trSchemeTheoretic{}闭包，并令 $\mu$ 为复合态射 $X'\to X\times_A\prod\overline{U}_i\to X$，于是我们有方图
\[
\begin{tikzcd}[column sep=large, row sep=large]
& X' \ar[dl, hook, "\text{闭嵌入}"'] \ar[dr, "\mu"] & \\
X\times_A\prod\overline{U}_i \ar[rr, "\text{射影}"'] \ar[d, "\text{固有}"'] && X \ar[d, "\text{固有}"] \\
\prod\overline{U}_i \ar[rr, "\text{射影}"'] && \Spec A
\end{tikzcd}
\]
（其中正方形是笛卡尔的）。态射 $\mu$ 是\trProjective{}的（因为它是两个\trProjective\trMorphism{}的复合，而 $X$ 是\trQuasicompact{}的，\mbox{练习~17.3.I}）。我们将这样结束论证：证明 $\mu^{-1}(U)=U$（更确切地说，$\mu$ 在 $U$ 上是同构），以及 $X'\to\prod\overline{U}_i$ 是闭嵌入（由此，复合
\[
X'\longrightarrow\prod\overline{U}_i\longrightarrow\Spec A
\]
是\trProjective{}的）。\footnote{原书此处复合式首项作『$X$』，当为『$X'$』之误，径改。——译者注}

\textbf{18.9.E.（练习）}\quad 设 $T_0,\dots,T_n$ 是 $A$ 上一些\emph{\trSeparated}的\trScheme，各带有互相同构的开集，我们在每种情形都马虎地把这些开集称作 $V$。那么 $V$ 是 $T_0\times_A\cdots\times_A T_n$ 的一个\trLocally\trClosed\trSubscheme。令 $\overline{V}$ 为这个\trLocally\trClosed\trSubscheme{}的闭包。证明：
\begin{align*}
V &\cong \overline{V}\cap(V\times_A T_1\times_A\cdots\times_A T_n)\\
 &= \overline{V}\cap(T_0\times_A V\times_A T_2\times_A\cdots\times_A T_n)\\
 &= \cdots\\
 &= \overline{V}\cap(T_0\times_A\cdots\times_A T_{n-1}\times_A V).
\end{align*}
（第一个同构的提示：态射 $V\to T_1\times_A\cdots\times_A T_n$ 的图像是一个闭嵌入，因为 $T_1\times_A\cdots\times_A T_n$ 在 $A$ 上是\trSeparated{}的，见\mbox{命题~11.2.6}。于是 $V$ 在 $V\times_A T_1\times_A\cdots\times_A T_n$ 中的\trSchemeTheoretic{}闭包就是 $V$ 自身。最后，\trSchemeTheoretic{}闭包可以局部地计算，本质上依据\mbox{定理~9.4.4}。）

\textbf{18.9.F.（练习）}\quad 用（前一道练习背后的想法）证明 $\mu^{-1}(U)=U$。

只须再证 $X'\to\prod\overline{U}_i$ 是闭嵌入。注意 $X'\to\prod\overline{U}_i$ 是\trClosed{}的（它是两个\trClosedMap{}的复合），因此只须证明 $X'\to\prod\overline{U}_i$ 是\trLocallyClosedEmbedding。

\textbf{18.9.G.（练习）}\quad 令 $A_i$ 为 $U$ 在
\[
B_i:=X\times_A\overline{U}_1\times_A\cdots\times_A U_i\times_A\cdots\times_A\overline{U}_n
\]
（只有第 $i$ 项缺上横线）中的闭包，并令 $C_i$ 为 $U$ 在
\[
D_i:=\overline{U}_1\times_A\cdots\times_A U_i\times_A\cdots\times_A\overline{U}_n.
\]
中的闭包。证明：存在由投射 $B_i\to D_i$ 诱导的同构 $A_i\xrightarrow{\sim}C_i$。提示：注意投射 $B_i\to D_i$ 的\trSection{} $D_i\to B_i$——非正式地写作 $(t_1,\dots,t_n)\mapsto(t_i,t_1,\dots,t_n)$——是一个闭嵌入，因为它可以解释为到一个\trSeparated{} $A$-\trScheme{}的某个态射的图像（\mbox{命题~11.2.6}）。于是 $U$ 可以解释为 $D_i$ 的一个\trLocally\trClosed\trSubscheme，而后者又可以解释为 $B_i$ 的一个\trClosedSubscheme。这样一来，$U$ 在 $D_i$ 中的闭包便可以与它在 $B_i$ 中的闭包等同。

由于诸 $U_i$ 覆盖 $X$，诸 $\mu^{-1}(U_i)$ 覆盖 $\overline{X}$。但 $\mu^{-1}(U_i)=A_i$（闭包可以局部地计算——$U$ 在 $B_i$ 中的闭包是 $B_i$ 与 $U$ 在 $X\times_A\overline{U}_1\times_A\cdots\times_A\overline{U}_n$ 中的闭包 $\overline{X}$ 之交）。\footnote{原书此处作『closure can be be computed locally』，当为『closure can be computed locally』之误，衍一『be』字，径删。——译者注}

于是，在每个 $U_i$ 上，我们得到闭嵌入 $A_i\hookrightarrow D_i$，从而 $X'\to\prod\overline{U}_i$ 是\trLocallyClosedEmbedding，正是我们所要的。\hfill$\square$

\textbf{18.9.5.\ \trChow{}引理的其他版本。}\quad 我们不会用到这些版本，但它们的证明与我们前面已经证明的内容类似。

\textbf{18.9.6.}\ \emph{\trhRemark。}\quad 首先注意：若 $X$ 是\trReduced{}的（相应地，\trIrreducible{}的、\trIntegral{}的），则 $X'$ 也可以取成\trReduced{}的（相应地，\trIrreducible{}的、\trIntegral{}的）。

\textbf{18.9.H.（练习）}\quad 在上面的证明中适当地划去若干行，把假设「$\pi$ 是固有的」减弱为「$\pi$ 是有限型且分离的」，其代价是把结论「$\pi\circ\mu$ 是射影的」减弱为「$\pi\circ\mu$ 是拟射影的」。

\textbf{18.9.7.}\ \emph{\trhRemark。}\quad 目标 $\Spec A$ 可以推广为一个满足下列条件之一的概形 $S$：$(i)$ \trNoetherian；或者 $(ii)$ \trSeparated{}且\trQuasicompact{}，并具有有限多个\trIrreducibleComponent{}。这可以与\mbox{注记~18.9.6} 及\mbox{练习~18.9.H} 相结合。证明见 [Gr-EGA, II.5.6.1]。另见 [GW, Thm. 13.100]，那里有一个稍微更一般的版本。
